% ============================================================
%  NOTATION & CONVENTIONS  (remove before submission)
% ============================================================
\section*{Notation \& Conventions}
\addcontentsline{toc}{section}{Notation \& Conventions}

This page collects all notational conventions used in this thesis.
It is to be \textbf{removed before submission}.

% ------------------------------------------------------------------
\subsection*{1. Manifolds, Maps, and the Differential}
% ------------------------------------------------------------------

\begin{longtable}{@{} p{4cm} p{10cm} @{}}
	\toprule
	\textbf{Symbol} & \textbf{Meaning \& Convention} \\
	\midrule
	\endhead
	
	$M, N$ & smooth manifolds (always assumed without boundary unless stated) \\[4pt]
	
	$T_p M,\; T^*_p M$ & tangent / cotangent space at $p \in M$ \\[4pt]
	
	$TM,\; T^*M$ & tangent / cotangent bundle \\[4pt]
	
	$\diff f$ & exterior differential of $f\colon M\to\R$; a $1$-form on $M$ \\[4pt]
	
	$\dat{p} f = (\diff f)_p$ & differential of $f$ \emph{at the point} $p$;
	a linear map $T_p M \to \R$.
	\newline
	\textbf{Rule:} write $\dat{p}f$ (macro \texttt{\textbackslash dat\{p\}f})
	when evaluating at a specific point.
	Write $\diff f$ for the full $1$-form / bundle map.\\[4pt]
	
	$\dat{p} F$ & for $F\colon M\to N$: the linear map $T_p M\to T_{F(p)}N$.
	Also written $({\diff F})_p$; never write $\frac{\diff}{\diff x}F$
	for a map between manifolds. \\[4pt]
	
	$\frac{\diff}{\diff t}$
	& total derivative w.r.t.\ a \emph{single} real variable $t$
	(roman~$\diff$, macro \texttt{\textbackslash diff}).
	Use for curves, flows evaluated on a trajectory. \\[4pt]
	
	$\frac{\partial}{\partial t}$
	& partial derivative; use only when \emph{at least two independent}
	variables/parameters are present.
	\textbf{Never} write $\partial/\partial t$ for the time-derivative along
	a flow line. \\[4pt]
	
	$\nabla_v f$ & covariant derivative of $f$ in direction $v$
	(Levi-Civita or chosen connection) \\[4pt]
	
	$\grad_g f$ & gradient of $f$ with respect to Riemannian metric $g$;
	characterised by $g(\grad_g f, -) = \diff f$ \\[4pt]
	
	\bottomrule
\end{longtable}

% ------------------------------------------------------------------
\subsection*{2. Flows and Dynamical Systems}
% ------------------------------------------------------------------

\begin{longtable}{@{} p{4cm} p{10cm} @{}}
	\toprule
	\textbf{Symbol} & \textbf{Meaning \& Convention} \\
	\midrule
	\endhead
	
	$\phi_t(x)$
	& flow of a vector field at time $t$ through initial point $x$.
	\textbf{Convention:} always write $\phi_t$, not $\phi(t,\cdot)$;
	think of $\phi_t\in\Diff(M)$ as a one-parameter family of
	diffeomorphisms.\\[4pt]
	
	$\phi_t = \mathrm{e}^{t X}$
	& flow of the vector field $X$; $\phi_0 = \id$. \\[4pt]
	
	$\frac{\diff}{\diff t}\phi_t(x)$
	& time-derivative of the flow line $t\mapsto\phi_t(x)$;
	equals $X(\phi_t(x))$. Use roman $\diff$, never $\partial$. \\[4pt]
	
	$\Wu{p},\;\Ws{p}$
	& unstable / stable manifold of a critical (rest) point $p$
	(macros \texttt{\textbackslash Wu\{p\}}, \texttt{\textbackslash Ws\{p\}}) \\[4pt]
	
	$\Mflow{p}{q}$
	& moduli space of (unparametrised) flow lines from $p$ to $q$
	(macro \texttt{\textbackslash Mflow\{p\}\{q\}}) \\[4pt]
	
	\bottomrule
\end{longtable}

% ------------------------------------------------------------------
\subsection*{3. Restrictions and Evaluations}
% ------------------------------------------------------------------

\medskip\noindent
\textbf{Core rule:} distinguish carefully between a \emph{fiber},
a \emph{restriction to a subspace}, and a \emph{value/evaluation}.

\begin{longtable}{@{} p{4.5cm} p{9.5cm} @{}}
	\toprule
	\textbf{Symbol} & \textbf{Meaning \& Convention} \\
	\midrule
	\endhead
	
	$E_p$
	& \textbf{fiber} of a vector bundle $E\to X$ over $p\in X$.
	\textbf{Preferred over} $E|_{\{p\}}$ for fibers;
	use macro \texttt{E\_p}. \\[4pt]
	
	$\restr{E}{A}$
	& \textbf{restriction} of the bundle $E$ to the subspace
	$A\subseteq X$; this is itself a bundle over $A$.
	Macro: \texttt{\textbackslash restr\{E\}\{A\}}.\\[4pt]
	
	$\restr{\xi}{A}$
	& section $\xi\in\Gamma(E)$ restricted to $A$;
	yields $\xi|_A\in\Gamma(\restr{E}{A})$. \\[4pt]
	
	$\xi(p)$ or $\xi_p$
	& \textbf{value} of the section $\xi$ at the point $p$;
	an element of $E_p$. \textbf{Avoid} $\xi|_p$ for this; use $\xi(p)$.\\[4pt]
	
	$\omega_p = \restr{\omega}{p}$
	& value of a differential $k$-form $\omega$ at $p$;
	an element of $\Lambda^k T^*_p M$.
	Acceptable to write $\omega|_p$ only when
	the point-evaluation is being \emph{emphasised} over the fiber. \\[4pt]
	
	\textbf{Summary}
	& $E_p$ = fiber (bundle); \quad
	$\restr{E}{A}$ = restricted bundle; \quad
	$\xi(p)$ = value of section; \quad
	$\omega_p$ = value of form. \\[4pt]
	
	\bottomrule
\end{longtable}

% ------------------------------------------------------------------
\subsection*{4. Exterior Derivative, Coboundary, Connecting Homomorphism}
% ------------------------------------------------------------------

\medskip\noindent
Three distinct operators share similar notation; use them as follows.

\begin{longtable}{@{} p{2cm} p{3cm} p{9cm} @{}}
	\toprule
	\textbf{Symbol} & \textbf{Macro} & \textbf{Use} \\
	\midrule
	\endhead
	
	$\partial$
	& \texttt{\textbackslash partial}
	& \textbf{Boundary operator} in chain complexes (homology direction).
	Default for singular chains; decorated for other theories
	(see §5). \\[4pt]
	
	$\diff$
	& \texttt{\textbackslash diff}
	& \textbf{Exterior derivative} on differential forms:
	$\diff\omega\in\Omega^{k+1}(M)$ for $\omega\in\Omega^k(M)$.
	Also the \emph{total derivative} w.r.t.\ one real variable
	(e.g.\ $\diff f/\diff t$). \textbf{Never} use $\diff$ as a
	boundary operator in a chain complex. \\[4pt]
	
	$\delta$
	& \texttt{\textbackslash delta}
	& \textbf{Two uses} (context distinguishes them):
	(1) \emph{Coboundary operator} in a cochain complex
	(cohomology direction), e.g.\ in singular cohomology
	$\delta\colon C^k\to C^{k+1}$.
	(2) \emph{Connecting homomorphism} in a long exact sequence
	of a pair or a short exact sequence of complexes,
	e.g.\ $\delta\colon H_n(X,A)\to H_{n-1}(A)$.
	Decorate when both appear simultaneously:
	$\delta^{\mathrm{conn}}$ vs.\ $\delta^{\mathrm{cob}}$.\\[4pt]
	
	\bottomrule
\end{longtable}

% ------------------------------------------------------------------
\subsection*{5. Chain Complexes and Homology Theories}
% ------------------------------------------------------------------

\medskip\noindent
Three homology theories appear simultaneously; they are always decorated.

\begin{longtable}{@{} p{3.5cm} p{3.5cm} p{7cm} @{}}
	\toprule
	\textbf{Theory} & \textbf{Chain complex} & \textbf{Boundary / Homology} \\
	\midrule
	\endhead
	
	Singular
	& $(\ChSing_*(X),\,\partial)$
	& boundary $\partial$; homology $H_*(X)$.
	\emph{Default / undecorated.} \\[4pt]
	
	CW
	& $(\ChCW_*(X),\,\bCW)$
	& boundary $\bCW$; homology $\HCW_*(X)$.
	Macros: \texttt{\textbackslash ChCW}, \texttt{\textbackslash bCW},
	\texttt{\textbackslash HCW}. \\[4pt]
	
	Morse
	& $(\ChM_*(f),\,\bM)$
	& boundary $\bM$; homology $\HM_*(f)$.
	Macros: \texttt{\textbackslash ChM}, \texttt{\textbackslash bM},
	\texttt{\textbackslash HM}. \\[4pt]
	
	\midrule
	\multicolumn{3}{@{}l}{%
		\textbf{Connecting homomorphisms in long exact sequences:} always $\delta$
		(decorated if needed: $\delta^{\mathrm{M}}, \delta^{\mathrm{CW}}$).}\\[4pt]
	
	\multicolumn{3}{@{}l}{%
		\textbf{Filtration:} $F^p\ChSing_*(X)$ for the filtration of the
		singular complex; boundary is still $\partial$.}\\[4pt]
	
	\bottomrule
\end{longtable}

% ------------------------------------------------------------------
\subsection*{6. Vector Bundles and K-Theory}
% ------------------------------------------------------------------

\begin{longtable}{@{} p{4cm} p{10cm} @{}}
	\toprule
	\textbf{Symbol} & \textbf{Meaning} \\
	\midrule
	\endhead
	
	$\Kvect{X}$
	& set of isomorphism classes of complex vector bundles over $X$
	(macro \texttt{\textbackslash Kvect\{X\}}) \\[4pt]
	
	$[E]$
	& stable isomorphism class / element in $\K(X)$ \\[4pt]
	
	$\K(X) = \KU(X)$
	& complex K-theory (default); macro \texttt{\textbackslash K},
	\texttt{\textbackslash KU} \\[4pt]
	
	$\KO(X)$
	& real K-theory; macro \texttt{\textbackslash KO} \\[4pt]
	
	$\tilde{\K}(X)$
	& reduced complex K-theory \\[4pt]
	
	$E\oplus F,\; E\otimes F$
	& Whitney sum / tensor product of bundles \\[4pt]
	
	$\restr{E}{A}$
	& restriction of bundle $E\to X$ to subspace $A$
	(same convention as §3) \\[4pt]
	
	$E_p$
	& fiber of $E$ over $p$ (same convention as §3) \\[4pt]
	
	\bottomrule
\end{longtable}

% ------------------------------------------------------------------
\subsection*{7. Spectral Sequences}
% ------------------------------------------------------------------

\begin{longtable}{@{} p{4cm} p{10cm} @{}}
	\toprule
	\textbf{Symbol} & \textbf{Meaning} \\
	\midrule
	\endhead
	
	$\Epq{r}{p,q}$
	& $(p,q)$-entry on the $r$-th page of a spectral sequence;
	macro \texttt{\textbackslash Epq\{r\}\{p,q\}} \\[4pt]
	
	$\drpq{r}$
	& differential on the $r$-th page;
	$\drpq{r}\colon \Epq{r}{p,q}\to \Epq{r}{p-r,\,q+r-1}$
	(cohomological convention);
	macro \texttt{\textbackslash drpq\{r\}} \\[4pt]
	
	AHSS
	& Atiyah--Hirzebruch spectral sequence:
	$\Epq{2}{p,q}\cong H^p(X;\,\K^q(\mathrm{pt}))
	\Rightarrow \K^{p+q}(X)$ \\[4pt]
	
	$\Epq{\infty}{p,q}$
	& abutment / $E_\infty$-page \\[4pt]
	
	\bottomrule
\end{longtable}

% ------------------------------------------------------------------
\subsection*{8. Frequently Used Macros (Quick Reference)}
% ------------------------------------------------------------------

\begin{longtable}{@{} p{4cm} p{5cm} p{5cm} @{}}
	\toprule
	\textbf{Macro} & \textbf{Output} & \textbf{Use} \\
	\midrule
	\endhead
	
	\texttt{\textbackslash diff}        & $\diff$           & exterior deriv.\ / total deriv. \\
	\texttt{\textbackslash dat\{p\}}    & $\dat{p}$         & deriv.\ at point $p$ \\
	\texttt{\textbackslash restr\{E\}\{A\}} & $\restr{E}{A}$ & restriction \\
	\texttt{\textbackslash Wu\{p\}}     & $\Wu{p}$          & unstable manifold \\
	\texttt{\textbackslash Ws\{p\}}     & $\Ws{p}$          & stable manifold \\
	\texttt{\textbackslash Mflow\{p\}\{q\}} & $\Mflow{p}{q}$ & moduli of flow lines \\
	\texttt{\textbackslash ChM}         & $\ChM_*$          & Morse chain complex \\
	\texttt{\textbackslash bM}          & $\bM$             & Morse boundary \\
	\texttt{\textbackslash HM}          & $\HM_*$           & Morse homology \\
	\texttt{\textbackslash ChCW}        & $\ChCW_*$         & CW chain complex \\
	\texttt{\textbackslash bCW}         & $\bCW$            & CW boundary \\
	\texttt{\textbackslash HCW}         & $\HCW_*$          & CW homology \\
	\texttt{\textbackslash conn}        & $\conn$           & connecting homom.\ ($\delta$) \\
	\texttt{\textbackslash K}           & $\K$              & K-theory \\
	\texttt{\textbackslash KO}          & $\KO$             & real K-theory \\
	\texttt{\textbackslash KU}          & $\KU$             & complex K-theory \\
	\texttt{\textbackslash Epq\{r\}\{p,q\}} & $\Epq{r}{p,q}$ & SS entry \\
	\texttt{\textbackslash drpq\{r\}}   & $\drpq{r}$        & SS differential \\
	\texttt{\textbackslash Crit}        & $\Crit$           & critical point set \\
	\texttt{\textbackslash ind}         & $\ind$            & Morse index \\
	\texttt{\textbackslash grad}        & $\grad$           & gradient \\
	\texttt{\textbackslash euler}       & $\euler$          & Euler number (roman e) \\
	\texttt{\textbackslash ii}          & $\ii$             & imaginary unit (roman i) \\
	\bottomrule
\end{longtable}